\chapter{Marco teórico y antecedentes}
\label{cap:marco}
\estadoBorrador

La revisión extensa de literatura continúa en desarrollo. Esta versión incorpora
la historia institucional de pensiones y transferencias; la introducción ya incorpora
referencias verificadas para motivar el problema; todavía se requiere una
revisión organizada sobre vulnerabilidad económica, patrimonio e ingreso,
salud y funcionamiento, predicción longitudinal y selección del seguimiento.
La revisión distinguirá estudios en México de evidencia de otros países,
diseños descriptivos de diseños causales y mediciones contemporáneas de
pronósticos. No se infiere una ausencia de literatura ni una novedad científica
definitiva a partir de la selección inicial de referencias.

% BEGIN CONSOLIDATED 02_recuperacion_conceptual
\section{Trayectoria laboral, protección pensionaria y posición económica}
\label{sec:recuperacion-conceptual}

Desde una perspectiva actuarial, el análisis de la seguridad económica en la
vejez no se limita al estudio de las reglas de financiamiento de las pensiones.
También requiere considerar la trayectoria laboral de las personas, las
condiciones de salud que pueden afectar su participación y la manera en que
distintos riesgos se acumulan a lo largo del ciclo de vida. Esta formulación,
recuperada del planteamiento original de la autora, delimita una perspectiva
de investigación: estudiar procesos que anteceden al deterioro económico,
sin asumir que cada retiro laboral o enfermedad conduce necesariamente a él.
La literatura sobre salud y recursos motiva esos posibles mecanismos, cuya
dirección y magnitud no se dan por demostradas en este panel \citep{smith1999}.

La situación económica de una persona durante la vejez puede relacionarse
con procesos que se desarrollan durante periodos prolongados. La educación,
las oportunidades laborales y la continuidad de la cotización son dimensiones
distintas de una trayectoria. Una pensión de supervivencia, por ejemplo,
corresponde al beneficiario que recibe la prestación, aunque el derecho
proceda de la historia contributiva de otra persona. La titularidad económica
no se resuelve mediante el parentesco ni por el hecho de que dos cónyuges
compartan hogar; los derechos derivados y las contingencias se distinguen en
las leyes respectivas \citep{lss1995,issste2007}.

\subsection{Regímenes de 1973 y 1997, cuentas individuales y AFORE}

En el IMSS coexisten pensiones asociadas al régimen de 1973 y al esquema
vigente desde julio de 1997. La nueva Ley del Seguro Social fue publicada en
1995; su fecha de publicación no debe confundirse con la entrada en vigor
del régimen habitualmente denominado de 1997. El acuerdo institucional de 2015 distingue a quienes cotizaron antes
del primero de julio de 1997. Las disposiciones transitorias de la ley
regulan la opción entre esquemas de quienes reúnen las condiciones previstas. Conocer que una persona declara IMSS en ENASEM no
establece cuál de esas opciones ejerció \citep{lss1995,imss2015regimenes}.

El sistema de cuentas individuales vincula la acumulación para el retiro con
una cuenta administrada por una AFORE. La Ley del Seguro Social regula la
cuenta y las modalidades de prestación; la identificación de una institución
otorgante no equivale a observar el saldo de esa cuenta, el historial de
cotización o la modalidad de pago. Para la tesis, esta distinción permite
separar la institución declarada de la historia laboral que sería necesaria
para explicar el acceso a determinada prestación \citep[arts.~159 y 175]{lss1995}.
No se asignan automáticamente los regímenes 1973/1997 según edad, institución
o año de entrevista. Tampoco se extrapolan a otras instituciones las reglas
del IMSS: la reforma del ISSSTE de 2007 tiene disposiciones y regímenes
transitorios propios \citep{issste2007}.

\subsection{Patrimonio, flujos de protección y vulnerabilidad}

El patrimonio neto ofrece una medida complementaria a los indicadores de
ingreso porque representa recursos acumulados y obligaciones, en lugar de
un flujo periódico. En ENASEMsimple, \variable{imp_neto} corresponde al
entrevistado y su cónyuge cuando procede, e incorpora imputación del
proveedor. Una misma cantidad patrimonial puede describir a una pareja
registrada en dos filas; ello no la convierte en dos patrimonios independientes
ni determina su reparto individual \citep{enasem2025simple}.

El borrador original justificaba la elección patrimonial porque la base
reducida no contenía ingresos ni pensiones directamente. Esa restricción
debe actualizarse: la integración provisional ya conserva las variables
económicas de diecisiete fuentes. Ahora es posible describir recepción,
instituciones y montos cuando su procedencia y titularidad se validan,
manteniendo el patrimonio como dimensión distinta. Esta ampliación no
resuelve automáticamente la definición de vulnerabilidad ni demuestra que
recibir una prestación impida una transición patrimonial adversa.

La vulnerabilidad puede formularse conceptualmente como susceptibilidad a
dificultades ante cambios laborales, familiares, de salud o de gasto. Una
posición patrimonial relativa es una posible aproximación empírica a esa
dimensión, cuya validez deberá justificarse. No se la equipara con pobreza
multidimensional, insuficiencia de ingreso ni una medida de incertidumbre de
bienestar \citep{ligon2003}. El stock patrimonial puede incluir activos de
liquidez limitada; su nivel y los flujos pensionarios describen aspectos
complementarios de la seguridad económica. La comprobación de liquidez,
necesidades de gasto y suficiencia permanece fuera de los resultados
descriptivos actualmente disponibles.

La aportación longitudinal consiste en distinguir el estado observado de
la entrada posterior a una condición de riesgo definida. Una variación de
posición relativa puede ocurrir por cambios en los recursos propios, en la
distribución de referencia o en la composición de la pareja. Por ello, usar
cuantiles por ronda no demuestra una pérdida real de poder adquisitivo.
La muestra que determina el corte, los precios, los empates y la unidad
económica deben documentarse antes de construir etiquetas. La formulación
provisional de estas decisiones aparece en la
sección~\ref{sec:recuperacion-metodologica}.
% END CONSOLIDATED 02_recuperacion_conceptual
% BEGIN CONSOLIDATED 02_pensiones_contexto
\section{Clasificación y evolución histórica de las pensiones contributivas, no contributivas y transferencias públicas en México (2001--2021).}
\label{sec:pensiones-contexto}

La protección económica durante la vejez reúne mecanismos con fundamentos
institucionales diferentes. Una pensión contributiva deriva de la participación
en un sistema de aseguramiento u ocupacional y puede corresponder a retiro,
invalidez o un derecho de sobrevivencia. Por tanto, una pensión por viudez no
se clasifica automáticamente como no contributiva. Las leyes del Seguro Social
y del ISSSTE distinguen contingencias, condiciones y regímenes; la institución
reportada permite identificar parte del sistema, pero no reconstruir por sí sola
la trayectoria de cotización ni el régimen individual
\citep{lss1995,issste2007,imss2006}.

Las prestaciones públicas de vejez sin requisito contributivo se distinguen,
a su vez, de las transferencias productivas, educativas, asistenciales o en
especie. Esta distinción es relevante para la tesis: patrimonio e ingreso
representan dimensiones diferentes de la posición económica. Dos personas con
un patrimonio registrado semejante pueden tener flujos de trabajo, pensiones
y apoyos distintos. Esta posibilidad orienta una comparación descriptiva,
no constituye un hallazgo ya medido ni demuestra un efecto causal de los
programas. Tampoco se equipara la institución que otorga un apoyo con su
naturaleza contributiva: una institución pública puede administrar prestaciones
que no dependen de cotizaciones.

\subsection{Antecedentes locales y expansión federal}

La documentación del Distrito Federal sitúa en marzo de 2001 el inicio del
apoyo alimentario y de atención dirigido a residentes de 70 años o más.
La ley de pensión alimentaria se publicó en noviembre de 2003 y la ampliación
a 68 y 69 años se aplicó desde 2009. Las reglas de 2017 para la Ciudad de México
establecían residencia permanente de al menos tres años y una transferencia
mensual de 1,132.35 pesos, vinculada a media Unidad de Medida y Actualización
diaria. Ese importe corresponde a la versión de 2017: no se proyecta hacia
el origen ni se interpreta como poder adquisitivo constante
\citep{evalua2015,cdmx2017}.

En el ámbito federal, las reglas de septiembre de 2003 del Programa de Atención
a los Adultos Mayores contemplaban personas de 60 años o más en pobreza
alimentaria, en localidades rurales de hasta 2,500 habitantes con alta o muy alta
marginación. Esa versión documenta un apoyo de 700 pesos mensuales y restricciones
sobre otros apoyos similares. Su existencia es un antecedente pertinente, pero
no acredita que todas las versiones posteriores hayan mantenido ese importe
ni que se trate del mismo programa que 70 y Más \citep{rural2003}.

Las reglas de 2007 establecieron el programa dirigido a personas de 70 años o
más en zonas rurales, con un apoyo mensual de 500 pesos y entrega ordinaria
bimestral. La cobertura territorial cambió: se documentan localidades de hasta
2,500 habitantes en 2007, hasta 20,000 en 2008 y hasta 30,000 en 2009; las reglas
para 2012 incluyeron localidades de mayor tamaño. Esta expansión modifica el
universo elegible, de modo que una variación posterior en cobertura de la
encuesta no se podría interpretar únicamente como cambio de comportamiento
individual \citep{setentamas2007,coneval2008mayores,setentamas2009,setentamas2012}.

En 2013 el programa pasó a denominarse Pensión para Adultos Mayores y redujo
la edad mínima a 65 años. Las reglas iniciales exigían no recibir ingresos por
jubilación o pensión contributiva y dispusieron un apoyo de 525 pesos mensuales,
con aplicación del nuevo importe a partir del segundo bimestre. La versión
para 2014 introdujo el criterio de ingreso pensionario contributivo no superior
a 1,092 pesos mensuales y un apoyo de 580 pesos, que también se documenta en
las reglas para 2018. El criterio de 1,092 pesos no debe atribuirse a la regla
inicial de febrero de 2013 \citep{pam2013,pam2014,pam2018}.

La Pensión para el Bienestar de las Personas Adultas Mayores modificó nuevamente
el diseño en 2019. Las reglas contemplaron 68 años como edad general, 65 para
personas indígenas y la continuidad de personas de 65 a 67 años inscritas en
el padrón de 2018. El importe de esa versión fue de 1,275 pesos mensuales,
entregados bimestralmente. Sus condiciones permiten distinguir el programa de
la focalización contributiva precedente; no debe suponerse exclusividad entre
una pensión contributiva y una prestación no contributiva \citep{bienestar2019}.

Las reglas iniciales de 2021 establecieron 1,350 pesos mensuales y conservaron
condiciones diferenciadas de edad, incluidas personas indígenas y afromexicanas.
El acuerdo publicado el 7 de julio de 2021, vigente desde el día siguiente,
amplió la edad general a 65 años y estableció 1,550 pesos mensuales. La reforma
ocurre dentro del año: para relacionarla con ENASEM deben conocerse la fecha
efectiva de entrevista y el periodo al que refiere cada pregunta. Una recepción
referida al año anterior no demuestra exposición a la ampliación de julio de
2021 \citep{bienestar2021a,bienestar2021b}.

\subsection{Transferencias adicionales y programas locales}

La protección pública no se limita a las pensiones de vejez. Oportunidades
incluyó un componente para personas adultas mayores de 70 años en familias
beneficiarias, con un importe inicial de 250 pesos mensuales durante la segunda
mitad de 2006. Se distingue ese componente del conjunto del programa.
PROSPERA reunió intervenciones diferentes; su continuidad institucional no
convierte todas las transferencias familiares en pensiones
\citep{oportunidades2006,oportunidades2007,prospera2015}.

En el ámbito productivo, el decreto de PROCAMPO de 1994, el componente PROAGRO
de las reglas agrícolas para 2014 y los lineamientos de Producción para el
Bienestar de 2019 permiten documentar antecedentes y cambios de padrones.
Estas relaciones no prueban equivalencia de importes, requisitos o beneficiarios
individuales. De forma semejante, el decreto de creación de la Coordinación
Nacional de Becas para el Bienestar Benito Juárez de 2019 no permite homologar
indiscriminadamente todos los componentes de PROSPERA como una misma prestación
\citep{procampo1994,proagro2014,produccion2019,becas2019}.

El catálogo incorpora ejemplos locales pertinentes, además del Distrito
Federal. Las reglas de Jalisco para 2014 contemplaban personas de 65 años o más,
residencia de tres años y restricciones sobre pensiones o apoyos similares,
con un importe de 1,009.35 pesos mensuales. La ficha histórica de 2018 sirve
como otra versión, sin trasladarle automáticamente el importe de 2014.
El informe municipal de Mérida documenta el programa Llegó a Ti, de entrega
domiciliaria de despensas a personas adultas mayores enfermas que no podían
recogerlas. Se trata de apoyo en especie, no de una pensión monetaria. Cuando
no se cotejó la creación exacta de un programa local, su año de origen permanece
desconocido \citep{jalisco2014,jalisco2018,merida2009}.

\begingroup\small\singlespacing
\setlength{\tabcolsep}{3pt}
\begin{longtable}{>{\raggedright\arraybackslash}p{1.1cm}>{\raggedright\arraybackslash}p{3.2cm}>{\raggedright\arraybackslash}p{1.6cm}>{\raggedright\arraybackslash}p{3.5cm}>{\raggedright\arraybackslash}p{3.7cm}>{\raggedright\arraybackslash}p{1.6cm}}
\caption{Programas pensionarios y transferencias: versiones históricas seleccionadas}\label{tab:pensiones_historia}\\
\toprule \textbf{Inicio} & \textbf{Nombre oficial} & \textbf{Ámbito} & \textbf{Requisitos relevantes} & \textbf{Cambios históricos} & \textbf{Fuente} \\
\midrule\endfirsthead
\toprule \textbf{Inicio} & \textbf{Nombre oficial} & \textbf{Ámbito} & \textbf{Requisitos relevantes} & \textbf{Cambios históricos} & \textbf{Fuente} \\
\midrule\endhead
\bottomrule\endfoot
2001 & Apoyo Alimentario, Atención Médica y Medicamentos Gratuitos para Adultos Mayores de 70 Años Residentes en el Distrito Federal & DF & Edad: 70+. Residencia DF; confirmar requisitos del instrumento inicial & 2001, antecedente documentado. Inicio en marzo de 2001; ley publicada 18/11/2003; expansión a 68--69 en 2009 & Evalúa CDMX \citeyearpar{evalua2015} \\
2001 & Pensión Alimentaria para Personas Adultas Mayores de 68 años, residentes en la Ciudad de México & CDMX & Edad: 68+. 68+ y residencia permanente CDMX no menor de tres años; solicitud y verificación & Reglas 2017. Edad reducida desde 2009, según antecedentes; no confundir con el programa federal & Gaceta CDMX \citeyearpar{cdmx2017} \\
2003 & Programa de Atención a los Adultos Mayores & Federal & Edad: 60+. Pobreza alimentaria; localidades rurales hasta 2,500 habitantes, alta/muy alta marginación; sin apoyo similar & Reglas 2003. Denominación rural frecuente en antecedentes; mantener separado de 70 y Más. Versiones posteriores del monto pendientes & DOF \citeyearpar{rural2003} \\
2007 & Programa de Atención a los Adultos Mayores de 70 años y más en zonas rurales & Federal & Edad: 70+. Localidades hasta 2,500 habitantes; no apoyo federal similar; identidad, edad y domicilio & Reglas 2007. Inicio de 70 y Más; permite tránsito desde otros apoyos sin duplicación; no prueba identidad con el programa rural de 2003 & DOF \citeyearpar{setentamas2007} \\
2007 & Programa 70 y más & Federal & Edad: 70+. Localidades hasta 30,000 habitantes; revisar incompatibilidades de la versión & Reglas 2009. Expansión: hasta 20,000 en 2008 y hasta 30,000 en 2009; 2008 corroborado por evaluación institucional & DOF \citeyearpar{setentamas2009} \\
2007 & Programa 70 y Más & Federal & Edad: 70+. Cobertura nacional; sin ingresos por jubilación/pensión; incorporación priorizada y presupuestada & Reglas 2012. Incorpora localidades mayores de 30,000; conservar padrones anteriores y reglas del año & DOF \citeyearpar{setentamas2012} \\
2013 & Programa Pensión para Adultos Mayores & Federal & Edad: 65+. Cobertura nacional; no recibir ingreso por jubilación o pensión contributiva & Reglas 2013; denominación desde 2013. Abroga reglas 70 y Más 2012. No trasladar el umbral de 1,092 de versiones posteriores a la regla inicial de 2013 & DOF \citeyearpar{pam2013} \\
2013 & Programa Pensión para Adultos Mayores & Federal & Edad: 65+. No recibir ingreso mensual superior a 1,092 por jubilación o pensión contributiva; revisar expediente y padrón & Reglas 2014. Monto aumentado a 580; regla anual separada de 2013 & DOF \citeyearpar{pam2014} \\
2013 & Programa Pensión para Adultos Mayores & Federal & Edad: 65+. No recibir pensión contributiva superior a 1,092 mensuales; requisitos del año & Reglas 2018 antes del rediseño 2019. Última versión anual previa al rediseño; no extender criterios a 2019 & DOF \citeyearpar{pam2018} \\
2019 & Pensión para el Bienestar de las Personas Adultas Mayores & Federal & Edad: 68 general; 65 indígenas; continuidad 65--67 del padrón 2018. Residencia y documentación; reglas ya no excluyen universalmente por tener pensión contributiva & Rediseño desde 2019. Nuevo nombre y diseño. Transitorio noveno conserva personas 65--67 previamente inscritas; no asumir elegibilidad general 65 en 2019 & DOF \citeyearpar{bienestar2019} \\
2019 & Programa Pensión para el Bienestar de las Personas Adultas Mayores & Federal & Edad: 68 general; 65 indígenas/afromexicanos; continuidad padrón. Residencia mexicana; identificación y criterios diferenciados del texto inicial & 2021 antes de modificación julio. Reglas iniciales vigentes desde 01/01/2021; no representan todo el año & DOF \citeyearpar{bienestar2021a} \\
2019 & Programa Pensión para el Bienestar de las Personas Adultas Mayores & Federal & Edad: 65+. 65+ general; mexicanas por nacimiento/naturalización y domicilio en México & Modificación publicada 07/07/2021. Ampliación general a 65 y mayor monto; acuerdo entra en vigor 08/07/2021, no desde enero & DOF \citeyearpar{bienestar2021b} \\
2006 & Programa de Desarrollo Humano Oportunidades: componente para adultos mayores & Federal & Edad: 70+. Integrantes de familias beneficiarias; corresponsabilidades y restricciones del componente & Componente 2006. Componente específico, distinto de todo el ingreso familiar Oportunidades y de 70 y Más & DOF \citeyearpar{oportunidades2006} \\
2014 & PROSPERA Programa de Inclusión Social & Federal & Edad: Según componentes, no edad única. Familias en pobreza y corresponsabilidades específicas; no toda transferencia es pensión & Reglas 2015. Cambio institucional desde Oportunidades; mantener concepto de cada componente. ENASEM puede agrupar etiquetas & DOF \citeyearpar{prospera2015} \\
1994 & Programa de Apoyos Directos al Campo (PROCAMPO) & Federal & Edad: No umbral general de vejez. Productores y superficies elegibles; no inferir pensión por edad del receptor & Creación jurídica 1994; antecedentes fuera de ventana. 1994 es fecha del decreto; inicio operativo previo requiere fuente específica; no confundir subsidio productivo y pensión & DOF \citeyearpar{procampo1994} \\
2014 & PROAGRO Productivo, componente del Programa de Fomento a la Agricultura & Federal & Edad: No umbral general de vejez. Requisitos de productor, predio y componente; revisar cuotas según estrato & Ejercicio 2014. Rediseño de apoyo productivo; vinculación histórica con PROCAMPO no autoriza igualdad de montos o reglas & DOF \citeyearpar{proagro2014} \\
2019 & Programa Producción para el Bienestar & Federal & Edad: No umbral general de vejez. Pequeños/medianos productores y predios; incorporación de directorios PROAGRO, PIMAF y localidades indígenas & Lineamientos 2019. Integra antecedentes distintos; no simple cambio de nombre de un único programa & DOF \citeyearpar{produccion2019} \\
2019 & Becas para el Bienestar Benito Juárez (coordinación y programas específicos) & Federal & Edad: Elegibilidad educativa; no edad única de vejez. Distinguir programas educativos, estudiante beneficiario y hogar receptor & Decreto 2019. Reasignación desde PROSPERA y otros presupuestos; no equivalencia de todos sus componentes & DOF \citeyearpar{becas2019} \\
ND & Atención a los Adultos Mayores & Jalisco & Edad: 65+. Residencia Jalisco tres años; sin pensión de igual/superior monto; criterios económicos y presupuesto & Reglas 2014. Se documenta reducción 70$\rightarrow$65 en 2013, no inicio del programa en 2013. Separado del federal & Jalisco \citeyearpar{jalisco2014} \\
ND & Atención a los adultos mayores & Jalisco & Edad: 65+. 65+ sin ingreso por pensión contributiva; revisar requisitos completos de las reglas anuales & Ficha institucional 2018. Operación institucional corroborada en 2018; no extrapolar el monto de 2014 & Jalisco \citeyearpar{jalisco2018} \\
ND & Llego a Ti; entrega domiciliaria de despensas para adultos mayores & Mérida & Edad: Adultos mayores; mínimo exacto no documentado en informe. Beneficiarios impedidos por enfermedad para recoger apoyos; cotejar reglas originales & Segundo informe 2007--2010 (2009). Ejemplo municipal: transferencia en especie, no pensión contributiva ni ingreso monetario automáticamente & Mérida \citeyearpar{merida2009} \\
ND & Pensiones del Instituto Mexicano del Seguro Social & Federal & Edad: Por contingencia y régimen; no una edad universal. Cotización y derechos del asegurado o derivados; retiro, invalidez, riesgos y supervivencia & Regímenes coexistentes en 2001--2021. Ley publicada 1995, régimen usualmente denominado 1997; año de ley no es creación de todas las pensiones IMSS & DOF \citeyearpar{lss1995} \\
ND & Pensiones del Instituto de Seguridad y Servicios Sociales de los Trabajadores del Estado & Federal & Edad: Por modalidad y régimen; no una edad universal. Servicios/cotizaciones y derechos del trabajador o beneficiario; régimen transitorio frente a cuentas individuales & Reforma 2007 y regímenes de transición en ventana. Ley 2007 no es inicio de todas las pensiones ISSSTE; el nombre institucional no identifica régimen de cada persona & DOF \citeyearpar{issste2007} \\
\end{longtable}
\noindent\footnotesize Fuente: elaboración documental, catálogo reproducible. ND: origen no verificado, distinto de ausencia de operación. Versiones y montos son históricos; no identifican recepción individual en ENASEM.\par
\endgroup


La tabla~\ref{tab:pensiones_historia} presenta versiones documentales
seleccionadas, con encabezados simples y reglas horizontales. El catálogo
reproducible que la genera añade institución, importe, periodicidad, operación
documentada y localizador de fuente. Es una selección de programas principales
y antecedentes locales pertinentes, no un censo de las 32 entidades y todos
los municipios. Los años de operación no cotejados y los montos no verificados
permanecen abiertos; no se completa una serie anual mediante supuestos.

\subsection{Historia institucional e identificación en la encuesta}

La homologación conserva una familia federal documentada para 70 y Más,
Pensión para Adultos Mayores y Bienestar, junto con sus versiones de elegibilidad.
El programa rural de 2003, la pensión del Distrito Federal y los programas
locales permanecen separados. Una relación de antecedente, sustitución de
padrones o cambio institucional no equivale a identidad de medición. El catálogo
contiene nombres originales y homologados, ámbito, tipo, institución, inicio
documentado y fuente verificable; los campos individuales no se rellenan a
partir de las reglas de operación.

Conviene distinguir tres tiempos: el inicio oficial del programa, la primera
recepción declarada por la persona y la primera ronda en que la encuesta
permite identificarla como receptora. Esta última puede estar afectada por
censura izquierda, ausencia temporal o cambios de cuestionario. La falta de
un programa federal en un año tampoco demuestra ausencia de otros apoyos.
En consecuencia, el contexto histórico delimita hipótesis e interpretación,
pero no sustituye las respuestas de ENASEM. La identificación efectiva y las
restricciones de las ocho variables se establecen en la
sección~\ref{sec:pensiones-metodo}.
% END CONSOLIDATED 02_pensiones_contexto
% BEGIN CONSOLIDATED 02_catalogo_nacional
\section{Oferta pública nacional y alcance de su observación en ENASEM}
\label{sec:catalogo-nacional}

La protección económica debe situarse en una arquitectura que combina
instituciones contributivas y oferta federal, estatal y municipal. Limitar
los antecedentes al Distrito Federal impediría reconocer diferencias de
oferta entre entidades y apoyos que no se denominan pensiones. Los
inventarios oficiales permiten ampliar ese contexto, pero su unidad es la
intervención pública documentada: no contienen el historial de recepción
de las personas del panel. La existencia de una intervención tampoco
demuestra que todos los residentes elegibles la recibieron.

Para esta etapa se revisaron las hojas principales de los inventarios
estatales de 2010, 2011, 2012, 2013, 2014, 2016, 2018 y 2021, el histórico
federal y cinco cortes municipales por entidad. Se extrajeron
\catTotal{} registros documentales: \catEstatal{} estatales,
\catFederal{} federales y \catMunicipal{} municipales. Son ocurrencias por
ejercicio o periodo y territorio; un mismo programa puede aparecer en
varias versiones. Estos números no se suman para estimar personas
beneficiarias ni programas históricos únicos
\citep{conevalhistoricoestatal,conevalhistoricofederal,conevalhistoricmunicipal}.

\subsection{Cobertura territorial, periodos y clasificación documental}

El cuadro~\ref{tab:catalogo-cobertura} presenta las 32 entidades y tres
ejercicios estatales coincidentes con rondas del estudio, junto con el
corte municipal de 2020. Las cifras son conteos de intervenciones
inventariadas y no tasas de cobertura poblacional. En 2021, por ejemplo,
Ciudad de México y Jalisco registran 84 intervenciones estatales cada una,
mientras Nayarit registra nueve. Esa diferencia permite reconocer
heterogeneidad documental de oferta, pero no establece diferencias de
protección efectiva, presupuesto por persona o calidad de atención.
Dentro de los programas existen beneficios, requisitos y tamaños de
población objetivo diferentes.

% Generado: no editar cifras manualmente.
\begingroup\small
\setlength{\tabcolsep}{4pt}
\begin{longtable}{lrrrr}
\caption{Cobertura documental por entidad: intervenciones en inventarios oficiales seleccionados}\label{tab:catalogo-cobertura}\\
\toprule
\textbf{Entidad} & \textbf{2012} & \textbf{2018} & \textbf{2021} & \textbf{Municipal 2020} \\
\midrule\endfirsthead
\multicolumn{5}{l}{\textit{Continuación}}\\
\toprule
\textbf{Entidad} & \textbf{2012} & \textbf{2018} & \textbf{2021} & \textbf{Municipal 2020} \\
\midrule\endhead
\bottomrule\endfoot
Aguascalientes & 55 & 33 & 27 & 114 \\
Baja California & 90 & 79 & 29 & 89 \\
Baja California Sur & 67 & 24 & 23 & 88 \\
Campeche & 118 & 39 & 23 & 146 \\
Coahuila & 123 & 39 & 37 & 308 \\
Colima & 130 & 40 & 17 & 81 \\
Chiapas & 253 & 19 & 17 & 221 \\
Chihuahua & 74 & 19 & 31 & 921 \\
Ciudad de México & 236 & 67 & 84 & 191 \\
Durango & 111 & 29 & 22 & 70 \\
Guanajuato & 99 & 93 & 57 & 670 \\
Guerrero & 138 & 46 & 31 & 216 \\
Hidalgo & 165 & 22 & 40 & 465 \\
Jalisco & 61 & 46 & 84 & 356 \\
Estado de México & 232 & 61 & 43 & 656 \\
Michoacán & 76 & 45 & 16 & 1\,286 \\
Morelos & 80 & 51 & 14 & 43 \\
Nayarit & 210 & 24 & 9 & 198 \\
Nuevo León & 89 & 21 & 27 & 362 \\
Oaxaca & 91 & 85 & 40 & 591 \\
Puebla & 125 & 36 & 22 & 1\,141 \\
Querétaro & 27 & 38 & 18 & 109 \\
Quintana Roo & 93 & 43 & 28 & 19 \\
San Luis Potosí & 57 & 42 & 14 & 58 \\
Sinaloa & 141 & 74 & 25 & 46 \\
Sonora & 193 & 58 & 29 & 344 \\
Tabasco & 57 & 24 & 15 & 131 \\
Tamaulipas & 114 & 32 & 43 & 42 \\
Tlaxcala & 40 & 32 & 34 & 172 \\
Veracruz & 241 & 23 & 28 & 834 \\
Yucatán & 99 & 118 & 53 & 336 \\
Zacatecas & 103 & 42 & 58 & 45 \\
\end{longtable}
\noindent{\footnotesize Fuente: extracción reproducible de inventarios CONEVAL. Conteos de intervenciones, no beneficiarios ni programas históricos únicos. Estatal: ejercicio fiscal; municipal: corte documental 2020. No mide cobertura poblacional. Las reglas, importes y continuidad se validan por separado. Para 2001/2003 no se dispone de un inventario estatal equivalente; 2015 no se sustituye por 2014/2016.}
\endgroup


La serie estatal no es anual y homogénea desde 2001. La hoja de 2010 no
incluye filas de intervenciones para Aguascalientes y Tlaxcala; su ausencia
documental no acredita ausencia de programas. Para 2001 y 2003 se cuenta
con reglas seleccionadas, pero no con un inventario estatal equivalente;
el corte de 2015 no se reemplaza por los de 2014 o 2016. El inventario
federal se restringió a 2004--2021 del libro que también contiene años
posteriores. Los cortes municipales 2012--2013, 2014--2015 y 2015--2016
son periodos de recopilación documental, no evidencia automática de
operación continua en ambos años. El informe municipal advierte
limitaciones de disponibilidad de documentos y distingue fuentes como
presupuesto, cuenta pública, informe y plan de desarrollo
\citep{conevalmunicipal2020}.

La clasificación documental preserva apoyos económicos, alimentación,
discapacidad, asistencia y otras intervenciones sociales. Un apoyo en
especie, un servicio, un subsidio productivo y una pensión representan
recursos o protección distintos. La población objetivo categórica de un
inventario no acredita requisitos exactos de edad ni permite reconstruir
un importe individual a partir del presupuesto total. Se conservan los
campos desconocidos y los nombres originales; las reglas anuales,
montos, periodicidades e inicio o terminación se verifican por programa.
Los marcadores de revisión por vejez o discapacidad solo organizan el
catálogo; no validan elegibilidad ni contributividad.

Como ejemplo de un beneficio diferente de la pensión de vejez, las
reglas iniciales de discapacidad permanente para 2021 contemplaban
rangos de edad y territorios específicos, con 1\,350 pesos mensuales
pagaderos bimestralmente. La modificación publicada el 21 de julio
acotó la edad adulta hasta antes de cumplir 65 años y mantuvo ese
importe en su tabla. La fecha de firma, 5 de julio, es distinta de la
publicación y de la vigencia, 22 de julio. Esta prestación no se
identifica automáticamente con una limitación funcional ni con la
pregunta K sobre invalidez \citep{discapacidad2021a,discapacidad2021b}.

\subsection{Correspondencia con la encuesta y lectura de los resultados}

El cuadro~\ref{tab:catalogo-instrumentos} distingue preguntas agregadas
y opciones nominativas. Antes de 2018, el bloque de donativos públicos
analizado no identifica el nombre del programa. En 2018 aparecen
PROAGRO, PROSPERA e INAPAM; en 2021 cambian las opciones y se incorpora
explícitamente el programa de vejez. La categoría de otra institución
permanece abierta. Una fuente IMSS o ISSSTE en el bloque pensionario
identifica una institución, sin reconstruir régimen, cotizaciones o
expediente de beneficiario. Las prestaciones estatales y municipales
catalogadas carecen de una identificación nominativa general en estos
campos \citep{inegi2018k,inegi2021k}.

% Generado: no editar cifras manualmente.
\begin{table}[!htbp]
\centering
\caption{Especificidad del bloque de donativos públicos por ronda}\label{tab:catalogo-instrumentos}
\small\setlength{\tabcolsep}{4pt}
\renewcommand{\arraystretch}{1.15}
\begin{tabular}{lp{3.3cm}p{8cm}}
\toprule
\textbf{Año} & \shortstack[c]{\textbf{Recepción}\\\textbf{A propia}} & \shortstack[c]{\textbf{Identificación}\\\textbf{de programa}} \\
\midrule
2001 & k76a & sin programa nombrado \\
2003 & k79a & sin programa nombrado \\
2012 & k79a\_12 & sin programa nombrado \\
2015 & k79a\_15 & sin programa nombrado \\
2018 & k79a\_18 & 1: PROAGRO Productivo; 2: PROSPERA; 3: INAPAM; 4: Otra institución: tipo no resuelto \\
2021 & k79a\_21 & 1: Producción para el Bienestar; 2: Becas para el Bienestar Benito Juárez; 3: INAPAM; 4: Otra institución: tipo no resuelto; 5: Programa para el Bienestar de las Personas Adultas Mayores \\
\bottomrule
\end{tabular}
\par\medskip
\parbox{\linewidth}{\footnotesize Fuente: contratos documentados de cuestionarios y libros ENASEM; no se recalculan respuestas. Propio refiere a la rama económica del beneficiario validada, no necesariamente al informante. Cónyuge, códigos, slots y montos se conservan en la matriz completa. Antes de2018 no se observa nombre de programa en este bloque. La categoría otro y las fuentes públicas K permanecen indeterminadas.}
\end{table}


La correspondencia se establece entre documento, pregunta, categoría y
periodo; la asignación a una persona requiere respuesta y titularidad
validadas. El catálogo no se une a las personas mediante edad, residencia,
año o coincidencia de monto. PROSPERA no se equipara a IMSS-PROSPERA,
INAPAM no prueba una pensión monetaria y la familia de becas no identifica
automáticamente el componente educativo. Las transferencias adicionales
en dinero o especie son más amplias que las pensiones no contributivas.
Dos versiones documentales o dos respuestas de la pareja no representan
por sí mismas dos ingresos diferentes.

El EDA de protección económica del capítulo~\ref{cap:resultados} permite
examinar recepción dentro de perfiles de sexo, escolaridad, edad y
tamaño de localidad. Sus diferencias pueden ser compatibles con cambios
documentados de reglas e instrumentos, pero también con reclutamiento,
supervivencia y observación desconocida. No se identifica aquí cuánto
corresponde a cada mecanismo. La mayor frecuencia de apoyo en perfiles
de menor escolaridad o localidades de menos de cien mil habitantes
describe la muestra y no demuestra focalización efectiva ni un impacto.
Asimismo, el corte de localidad no identifica ruralidad estricta ni
permite estimar diferencias estatales de beneficiarios.

El catálogo ofrece contexto y restricciones de interpretación; los
resultados numéricos de recepción conservan sus propios denominadores
y permanecen sin ponderar. Los importes normativos son nominales y
dependen de la versión del programa. Compararlos como poder adquisitivo
constante o atribuir una reforma a una respuesta sobre el año anterior
exigiría verificaciones temporales y monetarias adicionales. Esta
separación mantiene abiertas las hipótesis sobre ampliación de cobertura
y superposición de prestaciones, sin presentarlas como conclusiones
causales ni convertir un catálogo parcial en un censo exhaustivo.
% END CONSOLIDATED 02_catalogo_nacional
